\chapter{Detailed Implementation}\label{chap:detailed-implementation}

This chapter describes the design and implementation of the Zero-Knowledge Logistic Regression (ZKLR) system. It covers the overall architecture, dataset preparation, fixed-point encoding, circuit design, proof-generation workflow, performance optimizations, and final results.

% ===========================================================================
\section{System Architecture}\label{sec:system-architecture}
% ===========================================================================

\subsection{Overview}
The ZKLR prototype follows a client–server architecture where the server holds a private logistic regression model and the client verifies predictions without learning the model parameters. The system consists of:

\begin{itemize}
    \item Python — dataset generation, model training, metric plotting
    \item Go — ZK circuit implementation, witness generation, proof generation, verification
    \item gnark — PLONK-based ZKP framework
\end{itemize}

\subsection{Technology Stack}
\begin{itemize}
    \item Go 1.24.1 with gnark v0.11.0
    \item BN254 curve with KZG commitments
    \item Python 3.8+ with pandas, NumPy, scikit-learn
\end{itemize}

% ===========================================================================
\section{Dataset Handling}\label{sec:dataset}
% ===========================================================================

\subsection{Dataset}
A synthetic dataset of 3,000 samples is generated with a single input feature:

\begin{itemize}
    \item Feature: student marks (\( x \in [0,100] \))
    \item Label rule: pass/fail threshold at 60
    \item Noise: 150 samples randomly flipped
\end{itemize}

\subsection{Dataset Generation (Python)}
Minimal example:

\begin{verbatim}
marks = np.random.randint(0, 101, size=3000)
labels = (marks >= 60).astype(int)
\end{verbatim}

\subsection{Model Training}
A logistic regression model is trained in Python and the parameters are exported:

\[
z = w x + b
\]
\[
\hat{y} = \sigma(z) = \frac{1}{1 + e^{-z}}
\]

Trained values:
\begin{itemize}
    \item Weight: \( w = -0.85735312 \)
    \item Bias: \( b = 50.94705066 \)
\end{itemize}

These are encoded into fixed-point format before circuit usage.

% ===========================================================================
\section{Fixed-Point Encoding}\label{sec:fixed}
% ===========================================================================

\subsection{Motivation}
ZKP circuits cannot use floating-point arithmetic. Instead, fixed-point encodings are used:

\begin{table}[h]
\centering
\begin{tabular}{|l|c|c|}
\hline
Component & Format & Precision \\
\hline
Linear computation & Q32 & \(2^{32}\) \\
Sigmoid input & Q10 & \(2^{10}\) \\
Sigmoid output & Q16 & \(2^{16}\) \\
\hline
\end{tabular}
\caption{Fixed-point formats used in ZKLR}
\end{table}

Example:

\begin{verbatim}
Q32_value = float_value * 2^32
\end{verbatim}

% ===========================================================================
\section{Circuit Design}\label{sec:circuit}
% ===========================================================================

\subsection{Circuit Structure}
The circuit enforces that the output prediction provided by the server is correct. It contains three logical parts:

\begin{enumerate}
    \item Pre-computed linear term \( z \) accepted as private witness
    \item Sigmoid approximation using lookup tables (8192 entries)
    \item Output equality check \(Y == \sigma(z)\)
\end{enumerate}

\subsection{Sigmoid via Lookup Table}
The table stores:

\[
\sigma(x) = \frac{1}{1 + e^{-x}}  
\]

for the range \([0,8]\), since values beyond this saturate.

Negative inputs are handled using the identity:

\[
\sigma(-z) = 1 - \sigma(|z|)
\]

\subsection{Constraint Count}
\begin{table}[h]
\centering
\begin{tabular}{|l|c|}
\hline
Component & Constraints \\
\hline
Sigmoid lookup table & 58,025 \\
Negative handling & 2 \\
Final output check & 1 \\
\hline
Total & \textbf{58,028} \\
\hline
\end{tabular}
\caption{Constraint distribution in the ZKLR circuit}
\end{table}

% ===========================================================================
\section{Proof Generation Workflow}\label{sec:workflow}
% ===========================================================================

\subsection{Setup (One-Time)}
\begin{itemize}
    \item Compile circuit
    \item Generate proving and verifying keys
\end{itemize}

\subsection{Per-Sample Proof Workflow}
For each sample:
\begin{enumerate}
    \item Compute fixed-point values
    \item Fill witness
    \item Generate proof
    \item Verify proof
\end{enumerate}

Both operations take constant time since proof size is fixed (896 bytes).

% ===========================================================================
\section{Performance Optimizations}\label{sec:optim}
% ===========================================================================

\subsection{Lookup Table Range Reduction}
Optimized from 30K entries to 8K entries by saturating beyond ±8.

\[
\sigma(8) \approx 0.9997, \quad \sigma(-8) \approx 0.0003
\]

This reduced constraints by 66\% and improved proving speed by 4.4×.

% ===========================================================================
\section{Results}\label{sec:results}
% ===========================================================================

\subsection{Performance Summary}
\begin{table}[h]
\centering
\begin{tabular}{|l|r|}
\hline
Metric & Value \\
\hline
Constraints & 58,028 \\
Proof size & 896 bytes \\
Avg prove time & 620 ms/sample \\
Avg verification time & 2 ms \\
Total time (3000 samples) & ~31 minutes \\
Accuracy & 97.2\% \\
\hline
\end{tabular}
\caption{Final performance measurements}
\end{table}

\subsection{Where to Attach Figures}
Attach your screenshots/plots here:

\begin{itemize}
    \item Figure: Proof size comparison  
    \item Figure: Proof time vs dataset size  
    \item Figure: Verification time  
    \item Figure: Setup time  
    \item Figure: Constraint count  
\end{itemize}

Suggested placements:

```latex
\begin{figure}[H]
\centering
\includegraphics[width=0.8\textwidth]{metrics_output/1_proof_size.png}
\caption{Proof size comparison across samples}
\end{figure}
